\chapter{Haar Measure and Weyl Integration}

\begin{center}\itshape
What does it mean to choose a rotation, or a unitary matrix, at random? There is exactly one fair way, the Haar measure, and on a compact matrix group it
comes from a volume form that the group carries along. Hermann Weyl turned integrals over $U(n)$ into integrals over the eigenvalues alone, and the
eigenvalues turn out to repel each other with a precise law. That law reappeared, to everyone's surprise, among the zeros of the Riemann zeta function.
\end{center}

On 6 April 1972, in the tea room of the Institute for Advanced Study in Princeton, the number theorist Sarvadaman Chowla asked the young Hugh Montgomery
whether he had met Freeman Dyson. Montgomery, a young number theorist visiting the Institute, had just found something curious about the zeros of the
Riemann zeta function, and by his own account he answered that he did not need to meet Dyson. Chowla dragged him across the
room anyway. Dyson asked politely what he was working on, and Montgomery said that he had been looking at how the zeros of the zeta function are
spaced, and that the density of pairs of zeros at distance $u$ (in the natural unit) seemed to be
\[1-\Bigl(\frac{\sin\pi u}{\pi u}\Bigr)^2.\]
Dyson answered at once that this was the pair correlation of the eigenvalues of a random Hermitian matrix, a formula he and others had computed ten
years earlier for the energy levels of heavy nuclei. Montgomery, as he told Karl Sabbagh decades later, had never heard of it, had never seen a random matrix,
and had hardly seen one since. The conversation lasted about five minutes, and it changed what mathematicians think the zeta function is related to.

The formula has a simple meaning. If points were scattered independently, the density of pairs at distance $u$ would be $1$ for all $u$; here it vanishes
at $u=0$, like $\frac{\pi^2}3u^2$. The zeros of the zeta function, like the eigenvalues of a random matrix, avoid each other. Behind Dyson's formula is a
theorem of Weyl from 1925, which says how the eigenvalues of a random unitary matrix are distributed, and behind Weyl's theorem is the measure that
makes ``random unitary matrix'' meaningful: the Haar measure. This chapter builds both, proves Dyson's formula for the unitary group, and so prepares the
bridge that Volume~11 will cross from matrices to primes.

Invariant integration on a group is older than Haar. In 1897 Adolf Hurwitz, wanting to average polynomials over all rotations to produce invariants,
wrote down an invariant volume on the orthogonal and unitary groups using Euler-angle coordinates; Weyl used it in 1925 to build the representation
theory of compact groups; Haar's theorem of 1933 (Volume~7, Chapter~8) made it a general principle. For matrix groups we can follow Hurwitz and simply
write the volume down.

\section{Invariant volume forms}

\emph{The measure of a volume form.} Let $M$ be an oriented manifold of dimension $m$ and $\omega$ an $m$-form on $M$. In an oriented chart, $\omega=a\,\dd x^1\wedge\dots\wedge\dd x^m$,
and we let $|\omega|$ be the Borel measure given in the chart by $|a|\dd\lambda$. The computation in the proof of Proposition~\ref{prop:intwelldef}, with $\det DT$ replaced
by $|\det DT|$, shows that these local measures agree on overlaps; a partition of unity glues them into one measure $|\omega|$ on $M$ (the gluing is the
same as in the definition of $\int_M\omega$, with $\mathbf 1_E\rho_j$ in place of $\rho_j\omega$). Three properties follow from the change of variables formula
(Volume~6, Theorem~8.8) in charts:
\begin{enumerate}
\item[(i)] if $\omega$ is positive in every oriented chart, then $\int f\dd|\omega|=\int_Mf\omega$ for continuous $f$ with compact support;
\item[(ii)] if $F\colon N\to M$ is a diffeomorphism, then $\int_N(g\circ F)\dd|F^*\omega|=\int_Mg\dd|\omega|$ for every Borel $g\ge0$ (orientation plays no role);
\item[(iii)] on a product $M\times N$, the measure $|\alpha\wedge\beta|$ of the product of top forms is the product measure $|\alpha|\otimes|\beta|$ (in product charts the density is $a(x)b(y)$),
so Fubini's theorem (Volume~6, Theorem~8.6) applies to it.
\end{enumerate}

Let $G$ be a matrix Lie group of dimension $m$ with Lie algebra $\mathfrak g$. By Theorem~\ref{thm:closedsubgroup} it is a submanifold of $M_n$ with $T_IG=\mathfrak g$, and since left
multiplication $L_g\colon x\mapsto gx$ is a diffeomorphism of $G$ with differential $X\mapsto gX$, its tangent space at $g$ is $T_gG=g\mathfrak g$. Write $R_g\colon x\mapsto xg$ for right multiplication.

\begin{theorem}[Haar measure of a matrix Lie group]\label{thm:haarform}
Let $\omega_I$ be a non-zero alternating $m$-form on $\mathfrak g$. The formula
\[\omega_g(gX_1,\dots,gX_m)=\omega_I(X_1,\dots,X_m)\qquad(X_i\in\mathfrak g)\]
defines a smooth nowhere vanishing $m$-form on $G$ with $L_h^*\omega=\omega$ for all $h$, and $\mu=|\omega|$ is a left Haar measure on $G$. Moreover
\[R_h^*\omega=\det\bigl(\mathrm{Ad}(h^{-1})\bigr)\,\omega,\qquad\text{so}\qquad\mu(Eh)=\bigl|\det\mathrm{Ad}(h^{-1})\bigr|\,\mu(E),\]
where $\mathrm{Ad}(h)\colon\mathfrak g\to\mathfrak g$ is the linear map $X\mapsto hXh^{-1}$.
\end{theorem}

\begin{proof}
\emph{Smoothness.} Choose a basis $X_1,\dots,X_m$ of $\mathfrak g$ and linear forms $\ell_1,\dots,\ell_m$ on $M_n$ with $\ell_i(X_j)=\delta_{ij}$ (extend the dual basis from $\mathfrak g$ to $M_n$
in any way). On the open set $GL_n$ the $1$-forms $\theta^i_g(v)=\ell_i(g^{-1}v)$ are smooth, since $g\mapsto g^{-1}$ is. Restricted to $G$ they satisfy $\theta^i_g(gX_j)=\delta_{ij}$, so
$\omega=c\,\theta^1\wedge\dots\wedge\theta^m$ with $c=\omega_I(X_1,\dots,X_m)\neq0$; this is smooth, and at each $g$ it is non-zero on the basis $gX_1,\dots,gX_m$ of $T_gG$.

\emph{Left invariance.} $(L_h^*\omega)_g(gX_1,\dots)=\omega_{hg}(hgX_1,\dots)=\omega_I(X_1,\dots)=\omega_g(gX_1,\dots)$. By (ii), $\int(\mathbf 1_{hE}\circ L_h)\dd|\omega|=\int\mathbf 1_{hE}\dd|\omega|$, that is
$\mu(E)=\mu(hE)$. In a chart $\mu$ has a continuous positive density, so it is finite on compact sets and positive on non-empty open sets: $\mu$ is a left Haar measure.

\emph{Right translations.} By Proposition~\ref{prop:liealgebra}, $h^{-1}Xh\in\mathfrak g$ for $X\in\mathfrak g$. Then
\[(R_h^*\omega)_g(gX_1,\dots,gX_m)=\omega_{gh}(gX_1h,\dots)=\omega_{gh}\bigl(gh\,h^{-1}X_1h,\dots\bigr)=\omega_I\bigl(\mathrm{Ad}(h^{-1})X_1,\dots\bigr),\]
which is $\det\mathrm{Ad}(h^{-1})\,\omega_I(X_1,\dots,X_m)$, an alternating form evaluated on the image of a linear map. Finally (ii) with $F=R_h$ and $g=\mathbf 1_{Eh}$ gives
$|\det\mathrm{Ad}(h^{-1})|\mu(E)=\mu(Eh)$.
\end{proof}

The function $\Delta(h)=|\det\mathrm{Ad}(h^{-1})|$ is the modular function of Volume~7, Chapter~8, now computed. For $GL_n(\R)$ it is $1$ (conjugation preserves volume on $M_n$, since
$\det(X\mapsto h^{-1}Xh)=\det(h)^{-n}\det(h)^n$), which agrees with Volume~7, Proposition~8.2; for the group of maps $x\mapsto ax+b$ it is not (Exercise~\ref{ex:axbmodular}).

\begin{corollary}\label{cor:compactunimodular}
If $G$ is compact, then $\Delta=1$: the measure $\mu$ is also right invariant. If $G$ is compact and connected, $R_h^*\omega=\omega$ for every $h$.
\end{corollary}

\begin{proof}
$h\mapsto\det\mathrm{Ad}(h^{-1})$ is a continuous homomorphism from $G$ to the multiplicative group $\R^\times$. Its image is a compact subgroup of $\R^\times$, and the only ones are $\{1\}$ and $\{\pm1\}$
(a number $x$ with $|x|\neq1$ has unbounded powers $x^k$ or $x^{-k}$). If $G$ is connected, the image is connected and contains $1$, so it is $\{1\}$.
\end{proof}

On a compact group we normalize: $\mu(G)=1$. The following theorem is the non-abelian, compact case of the uniqueness theorem of Volume~7 (Theorem~8.3), and its
proof is shorter, because compactness gives us a measure invariant on both sides to work with.

\begin{theorem}[uniqueness]\label{thm:haarcompact}
Let $G$ be a compact metrizable group carrying a right invariant Borel probability $\mu$. Every left invariant finite Borel measure $\nu$ on $G$ equals $\nu(G)\mu$.
In particular, on a compact matrix Lie group there is exactly one left invariant Borel probability, the Haar measure $\mu$ of Theorem~\ref{thm:haarform}; it is right
invariant, and invariant under inversion: $\mu(E^{-1})=\mu(E)$.
\end{theorem}

\begin{proof}
Let $f$ be continuous on $G$. The function $(x,y)\mapsto f(yx)$ is continuous on $G\times G$, so Fubini's theorem applies to $\nu\otimes\mu$. By left invariance of $\nu$,
$\int f(yx)\dd\nu(x)=\int f\dd\nu$ for every $y$; by right invariance of $\mu$, $\int f(yx)\dd\mu(y)=\int f\dd\mu$ for every $x$. Hence
\[\int f\dd\nu=\int_G\!\int_Gf(yx)\dd\nu(x)\dd\mu(y)=\int_G\!\int_Gf(yx)\dd\mu(y)\dd\nu(x)=\nu(G)\int f\dd\mu.\]
Two finite Borel measures with the same integrals of continuous functions on a compact metric space are equal: the end of the proof of Volume~7, Theorem~8.3
shows it, approximating indicators of compact sets by continuous functions.

For a compact matrix Lie group, $\mu$ is right invariant by Corollary~\ref{cor:compactunimodular}, so the first part applies. The measure $\check\mu(E)=\mu(E^{-1})$ is a left invariant
probability, since $(hE)^{-1}=E^{-1}h^{-1}$ and $\mu$ is right invariant; hence $\check\mu=\mu$.
\end{proof}

\begin{corollary}\label{cor:pushhaar}
Let $\varphi\colon G\to H$ be a continuous surjective homomorphism of compact matrix Lie groups. The image of the Haar measure of $G$ is the Haar measure of $H$.
\end{corollary}

\begin{proof}
Let $\nu(E)=\mu(\varphi^{-1}(E))$, a Borel probability on $H$. For $h=\varphi(g)$, $\varphi^{-1}(hE)=g\varphi^{-1}(E)$, so $\nu(hE)=\nu(E)$; conclude by Theorem~\ref{thm:haarcompact}.
\end{proof}

\begin{example}[$SU(2)$ and $SO(3)$]\label{ex:su2haar}
A matrix of $SU(2)$ is $\begin{pmatrix}a&-\bar b\\b&\bar a\end{pmatrix}$ with $|a|^2+|b|^2=1$, and is determined by its first column; this identifies $SU(2)$ with the unit sphere
$S^3\subset\C^2=\R^4$. Under this identification, left multiplication by $g\in SU(2)$ is the map $z\mapsto gz$ on first columns, which is unitary, hence an orthogonal map of $\R^4$.
The uniform probability $\sigma_3$ on $S^3$ is the law of $Z/|Z|$ for a standard Gaussian vector $Z$ in $\R^4$ (polar coordinates, as in the solution of Exercise~\ref{ex:spherevolume}), and
the Gaussian law is invariant under orthogonal maps. So $\sigma_3$ is a left invariant probability on $SU(2)$: \emph{the Haar measure of $SU(2)$ is the uniform measure on $S^3$.}
By Corollary~\ref{cor:pushhaar} and the double cover $SU(2)\to SO(3)$ of Theorem~\ref{thm:spin}, a Haar random rotation is obtained by taking a uniform point of $S^3$, a unit
quaternion, and the rotation it induces. This is how computer graphics draws random orientations (Ken Shoemake's method of 1992), and Exercise~\ref{ex:su2weyl} computes
how large the angle of a random rotation is.
\end{example}

\section{Sampling a random unitary matrix}

How does one produce a Haar random element of $U(n)$ on a computer? Take a matrix of independent Gaussian entries and orthonormalize its columns.
A \emph{standard complex Gaussian} is $\frac{X+iY}{\sqrt2}$ with $X,Y$ independent standard normal; its density on $\C=\R^2$ is $\frac1\pi e^{-|z|^2}$. A vector $z\in\C^n$ of independent
standard complex Gaussians has density $\pi^{-n}e^{-|z|^2}$, which is invariant under every unitary map $z\mapsto gz$; since $g$ is an orthogonal map of $\R^{2n}$, it preserves
Lebesgue measure too, so \emph{$gz$ has the same law as $z$}.

\begin{lemma}\label{lem:qr}
Every invertible $Z\in GL_n(\C)$ is uniquely $Z=QR$ with $Q$ unitary and $R$ upper triangular with positive real diagonal entries, and $Q$ depends continuously on $Z$. A random
matrix $Z$ with independent standard complex Gaussian entries is invertible with probability $1$.
\end{lemma}

\begin{proof}
Gram--Schmidt orthonormalization of the columns $z_1,\dots,z_n$ produces orthonormal $q_1,\dots,q_n$ with $z_k=\sum_{j\le k}r_{jk}q_j$ and $r_{kk}=|z_k-\sum_{j<k}\langle q_j,z_k\rangle q_j|>0$; each $q_k$
is a continuous function of $Z$, given by explicit formulas. If $QR=Q'R'$, then $Q'^{-1}Q=R'R^{-1}$ is unitary and upper triangular with positive diagonal; a unitary upper triangular
matrix is diagonal (its first column has norm $1$ and only one non-zero entry, so the other columns, orthogonal to it, vanish in the first row, and so on), with entries of modulus $1$,
here positive, so it is $I$. For the last claim, the columns are independent; given $z_1,\dots,z_{k-1}$, the column $z_k$ lies in their span, a subspace of real codimension at least
$2$, with probability $0$, because a proper subspace has Lebesgue measure zero. By Fubini, $z_1,\dots,z_n$ are linearly independent almost surely.
\end{proof}

\begin{theorem}[Gaussian sampling of Haar measure]\label{thm:qrhaar}
Let $Z$ be an $n\times n$ matrix with independent standard complex Gaussian entries, and $Z=QR$ its decomposition of Lemma~\ref{lem:qr}. Then $Q$ is distributed according to the Haar measure of $U(n)$.
\end{theorem}

\begin{proof}
Let $g\in U(n)$. The columns of $gZ$ are $gz_k$, independent and with the same law as the $z_k$, so $gZ$ has the same law as $Z$. Now $gZ=(gQ)R$ is the decomposition of $gZ$, by
uniqueness, so $Q(gZ)=gQ(Z)$. Therefore $gQ$ has the same law as $Q$: the law of $Q$ is a left invariant Borel probability on $U(n)$, and Theorem~\ref{thm:haarcompact} identifies it.
\end{proof}

\begin{remark}
Numerical libraries compute $QR$ by Householder reflections and do not make the diagonal of $R$ positive. The $Q$ they return is then \emph{not} Haar distributed, and the error is
visible in the eigenvalues. The cure, popularized by Francesco Mezzadri in 2007, is to replace $Q$ by $Q\Lambda$, where $\Lambda$ is the diagonal matrix of the phases $r_{jj}/|r_{jj}|$;
numerical evaluation shows both. Taking real Gaussians and the same recipe gives the Haar measure of $O(n)$.
\end{remark}

\section{Weyl's integration formula}

A function $f$ on $U(n)$ is a \emph{class function} if $f(gUg^{-1})=f(U)$ for all $g$, $U$: it depends only on the eigenvalues. The trace, the determinant, the characteristic polynomial, the
number of eigenvalues in an arc are class functions. By Theorem~\ref{thm:torus}, every $U$ is conjugate to a diagonal matrix $t=\mathrm{diag}(t_1,\dots,t_n)$ of the torus $T$, unique up to
the order of the entries. Write $t_j=e^{i\theta_j}$ and
\[\Delta(t)=\prod_{1\le j<k\le n}(t_k-t_j)=\det\bigl(t_k^{\,j-1}\bigr)_{1\le j,k\le n}\]
(the Vandermonde determinant of Volume~2).

\begin{theorem}[Weyl, 1925]\label{thm:weyl}
For every bounded Borel class function $f$ on $U(n)$,
\[\int_{U(n)}f\dd\mu=\frac1{n!}\int_{[0,2\pi]^n}f\bigl(\mathrm{diag}(e^{i\theta_1},\dots,e^{i\theta_n})\bigr)\prod_{j<k}\bigl|e^{i\theta_j}-e^{i\theta_k}\bigr|^2\,\frac{\dd\theta_1\cdots\dd\theta_n}{(2\pi)^n}.\]
In words: the eigenvalues of a Haar random unitary matrix, put in random order, have joint density $\frac1{n!}|\Delta|^2$ with respect to the uniform measure on $T$.
\end{theorem}

The proof parametrizes $U=gtg^{-1}$ by the eigenvalues $t$ and by the eigenvectors, and computes the Jacobian. The eigenvectors are only defined up to phases, so we record them by a
Hermitian matrix with the same eigenvectors and fixed eigenvalues $1,2,\dots,n$. Let $\mathrm{Herm}_n$ be the real vector space of Hermitian matrices, of dimension $n^2$, with the inner product
$\langle A,B\rangle=\operatorname{tr}(AB)$, an orientation and its Euclidean volume form $\mathrm{vol}$. For $g\in U(n)$, conjugation $c_g(H)=gHg^{-1}$ is a linear isometry of $\mathrm{Herm}_n$; its
determinant is continuous in $g$, equal to $\pm1$, and $U(n)$ is connected, so $\det c_g=1$ and $c_g^*\mathrm{vol}=\mathrm{vol}$.

\begin{lemma}\label{lem:orbit}
Let $D=\mathrm{diag}(1,2,\dots,n)$ and $\mathcal O=\{gDg^{-1}:g\in U(n)\}$. Then $\mathcal O$ is a compact submanifold of $\mathrm{Herm}_n$ of dimension $n^2-n$, its tangent space at $D$ is the space
of Hermitian matrices with zero diagonal, and it carries a nowhere vanishing $(n^2-n)$-form $\sigma$ with $c_g^*\sigma=\sigma$ for all $g\in U(n)$.
\end{lemma}

\begin{proof}
Let $p\colon\mathrm{Herm}_n\to\R^n$ be the map
\[p(H)=(\operatorname{tr}H,\operatorname{tr}H^2,\dots,\operatorname{tr}H^n).\]
The power sums $p_k=\sum_i\lambda_i^k$ ($k\le n$) of the eigenvalues determine their elementary symmetric functions
$e_k$, by Newton's identities $ke_k=\sum_{i=1}^k(-1)^{i-1}e_{k-i}p_i$ (compare coefficients in $x\,E'(x)=E(x)\sum_{k\ge1}(-1)^{k-1}p_kx^k$, where $E(x)=\prod_i(1+\lambda_ix)=\sum_ke_kx^k$), hence the characteristic
polynomial. So $p(H)=p(D)$ exactly when $H$ has eigenvalues $1,\dots,n$, that is, by the spectral theorem, when $H\in\mathcal O$: $\mathcal O=p^{-1}(p(D))$, closed and bounded, hence compact.

The differential is
\[dp_H(K)=\bigl(k\operatorname{tr}(H^{k-1}K)\bigr)_{k\le n}.\]
At $H=gDg^{-1}$ and $K=gE_{jj}g^{-1}$ (with $E_{jj}$ the diagonal matrix unit) it gives $(kj^{k-1})_{k\le n}$; for
$j=1,\dots,n$ these vectors are independent (a Vandermonde matrix with distinct nodes, times $\mathrm{diag}(k)$), so $dp_H$ is onto. By the regular value theorem
(Theorem~\ref{thm:regularvalue}), $\mathcal O$ is a submanifold of dimension $n^2-n$ with $T_H\mathcal O=\ker dp_H$. At $D$, $K\in\ker dp_D$ means $\sum_jj^{k-1}K_{jj}=0$ for all $k$, that is
$K_{jj}=0$ for all $j$.

The gradient of $p_k$ is $N_k(H)=kH^{k-1}$: indeed $dp_k(K)=\langle kH^{k-1},K\rangle$. At $H\in\mathcal O$ the $N_k(H)$ are orthogonal to $\ker dp_H=T_H\mathcal O$, and independent (a relation $\sum a_kN_k=0$
would give $\sum a_k\,dp_k=0$, contradicting surjectivity). Define
\[\sigma_H(K_1,\dots,K_{n^2-n})=\mathrm{vol}\bigl(N_1(H),\dots,N_n(H),K_1,\dots,K_{n^2-n}\bigr)\qquad(K_i\in T_H\mathcal O).\]
It is alternating, smooth (polynomial in $H$), and non-zero on a basis of $T_H\mathcal O$, which together with the $N_k(H)$ forms a basis of $\mathrm{Herm}_n$. Finally $N_k(c_gH)=c_gN_k(H)$ and $c_g$ is linear
with $c_g^*\mathrm{vol}=\mathrm{vol}$, so $c_g^*\sigma=\sigma$.
\end{proof}

For $H\in\mathcal O$, the Lagrange polynomials $P_j(x)=\prod_{k\neq j}\frac{x-k}{j-k}$ give the projections onto the eigenlines: $P_j(gDg^{-1})=gP_j(D)g^{-1}=gE_{jj}g^{-1}$, since $P_j(i)=\delta_{ij}$.
Define the smooth map
\[\Phi\colon\mathcal O\times T\to U(n),\qquad\Phi(H,t)=\sum_{j=1}^nt_jP_j(H),\qquad\text{so that}\quad\Phi(gDg^{-1},t)=gtg^{-1}.\]
Both sides have dimension $n^2$. Let $\omega$ be the Haar form of $U(n)$ (Theorem~\ref{thm:haarform}), and $\dd\theta=\dd\theta_1\wedge\dots\wedge\dd\theta_n$ on $T$ (the differentials $\dd\theta_j$ are well defined on $T$,
although $\theta_j$ is defined only modulo $2\pi$); the measure $|\dd\theta|$ is Lebesgue measure in the coordinates $\theta\in[0,2\pi)^n$.

\begin{lemma}[the Jacobian]\label{lem:jacobian}
There is a constant $C\neq0$ with $\Phi^*\omega=C\,|\Delta(t)|^2\,\sigma\wedge\dd\theta$ on $\mathcal O\times T$.
\end{lemma}

\begin{proof}
Write $\Phi^*\omega=J\,\sigma\wedge\dd\theta$ with $J$ smooth. Since $P_j(c_gH)=c_gP_j(H)$, we have $\Phi\circ(c_g\times\mathrm{id})=c_g\circ\Phi$. On $U(n)$, $c_g=L_g\circ R_{g^{-1}}$ preserves $\omega$
(Corollary~\ref{cor:compactunimodular}; $U(n)$ is connected), and $c_g\times\mathrm{id}$ preserves $\sigma\wedge\dd\theta$ (Lemma~\ref{lem:orbit}). Pulling back by $c_g\times\mathrm{id}$ therefore gives $J(c_gH,t)=J(H,t)$, and
it suffices to compute $J(D,t)$.

A tangent vector at $(D,t)$ is a pair $(K,it\Theta)$ with $K$ Hermitian of zero diagonal and $\Theta$ real diagonal (so that $\dd\theta_j$ takes the value $\Theta_{jj}$). Let $X$ be the
matrix with $X_{jj}=0$ and $X_{jk}=\frac{K_{jk}}{k-j}$ for $j\neq k$; it is anti-Hermitian, so $X\in\mathfrak u(n)$, and $[X,D]_{jk}=(k-j)X_{jk}=K_{jk}$. The curve $s\mapsto(e^{sX}De^{-sX},te^{is\Theta})$
has velocity $(K,it\Theta)$ at $s=0$, and $\Phi$ maps it to $e^{sX}te^{is\Theta}e^{-sX}$, whose velocity is
\[Xt-tX+it\Theta=t\bigl(t^{-1}Xt-X+i\Theta\bigr).\]
So $d\Phi_{(D,t)}(K,it\Theta)=tA(K,\Theta)$ with $A(K,\Theta)=t^{-1}Xt-X+i\Theta\in\mathfrak u(n)$, and $\omega_t(tY_1,\dots)=\omega_I(Y_1,\dots)$. Use the coordinates $(K_{jk})_{j<k}\in\C^{n(n-1)/2}$, $\Theta\in\R^n$ on the
tangent space at $(D,t)$, and the same coordinates (entries above the diagonal, and the diagonal divided by $i$) on $\mathfrak u(n)$. Since $(t^{-1}Xt)_{jk}=t_j^{-1}X_{jk}t_k$, the map $A$
multiplies the coordinate $K_{jk}$ by the complex number $z_{jk}=\frac{t_k/t_j-1}{k-j}$ and leaves $\Theta$ alone. Multiplication by $z$ on $\C=\R^2$ has determinant $|z|^2$, so
$\det A=\prod_{j<k}\frac{|t_k-t_j|^2}{(k-j)^2}$ (recall $|t_j|=1$). Hence $J(D,t)=C|\Delta(t)|^2$, where $C$ is $\prod_{j<k}(k-j)^{-2}$ times the ratio of $\omega_I$ to $\sigma_D\wedge\dd\theta$ on the coordinate bases,
a non-zero number independent of $t$.
\end{proof}

Call $t\in T$ \emph{regular} if its entries are distinct, and $U\in U(n)$ regular if its eigenvalues are distinct; write $T_{\mathrm{reg}}$, $U(n)_{\mathrm{reg}}$, and $T_{\mathrm{sing}}$, $U(n)_{\mathrm{sing}}$ for the rest.

\begin{lemma}\label{lem:sheets}
\begin{enumerate}
\item[(a)] $\Phi^{-1}(U(n)_{\mathrm{reg}})=\mathcal O\times T_{\mathrm{reg}}$, and every regular $U$ has exactly $n!$ preimages.
\item[(b)] Let $F\colon X\to Y$ be a local diffeomorphism between manifolds such that every fibre has exactly $N$ points. Then every $y\in Y$ has an open neighbourhood $W$ such that $F^{-1}(W)$
is the disjoint union of $N$ open sets, each mapped diffeomorphically onto $W$.
\end{enumerate}
\end{lemma}

\begin{proof}
(a) If $\Phi(H,t)=U$ with $H=gDg^{-1}$, then $U=gtg^{-1}$, so the entries of $t$ are the eigenvalues of $U$ and $t$ is regular exactly when $U$ is. Let $U$ be regular, with eigenvalues $\lambda_1,\dots,\lambda_n$
and eigenlines $\ell_1,\dots,\ell_n$. A preimage $(H,t)$ puts the eigenvalue $j$ of $H$ on the line $\ell_{w(j)}$ and $t_j=\lambda_{w(j)}$, for a permutation $w$; conversely each permutation $w$
gives the preimage $H=\sum_jj\,\pi_{w(j)}$, $t=\mathrm{diag}(\lambda_{w(j)})$, where $\pi_i$ is the orthogonal projection onto $\ell_i$. Distinct $w$ give distinct $t$.

(b) Let $F^{-1}(y)=\{x_1,\dots,x_N\}$. Choose disjoint open $V_i\ni x_i$ mapped diffeomorphically onto open sets $W_i\ni y$ (inverse function theorem, Theorem~\ref{thm:inverse}). Some open $W\ni y$ inside
$\bigcap_iW_i$ satisfies $F^{-1}(W)\subset\bigcup_iV_i$: otherwise there are points $x'_k\notin\bigcup V_i$ with $F(x'_k)\to y$, and for $k$ large $F(x'_k)\in\bigcap W_i$ has a preimage in each $V_i$ and
one more, $x'_k$, which is $N+1$ preimages. Then $F^{-1}(W)$ is the disjoint union of the sets $V_i\cap F^{-1}(W)$, each mapped diffeomorphically onto $W$.
\end{proof}

\begin{proof}[Proof of Theorem~\ref{thm:weyl}]
Orient $\mathcal O\times T$ so that $C>0$ in Lemma~\ref{lem:jacobian}.

\emph{Step 1: the singular set is negligible.} By Lemma~\ref{lem:jacobian}, the critical points of $\Phi$ are the points of $\mathcal O\times T_{\mathrm{sing}}$, and $\Phi(\mathcal O\times T_{\mathrm{sing}})=U(n)_{\mathrm{sing}}$
(a singular $U$ is $gtg^{-1}$ with $t$ singular, and $\Phi(gDg^{-1},t)=U$). Covering both manifolds by countably many charts and applying Sard's theorem (Theorem~\ref{thm:sard}) in each pair of charts, the critical
values form a set that is null in every chart of $U(n)$; since $\mu$ has a continuous density in charts, $\mu(U(n)_{\mathrm{sing}})=0$. Likewise $T_{\mathrm{sing}}$, a finite union of the subtori $\{t_j=t_k\}$, is Lebesgue null.

\emph{Step 2: counting sheets.} On $\mathcal O\times T_{\mathrm{reg}}$, $\Phi^*\omega$ vanishes nowhere, so $\Phi$ is a local diffeomorphism onto the open set $U(n)_{\mathrm{reg}}$, with fibres of exactly $n!$ points
(Lemma~\ref{lem:sheets}). For a Borel set $E$ contained in a neighbourhood $W$ as in Lemma~\ref{lem:sheets}(b), property (ii) on each of the $n!$ sheets gives
$\int_{\mathcal O\times T_{\mathrm{reg}}}\mathbf 1_E\circ\Phi\dd|\Phi^*\omega|=n!\,\mu(E)$. Countably many such $W$ cover $U(n)_{\mathrm{reg}}$ (it is a countable union of compact sets, each covered by finitely
many $W$), so cutting a general Borel $E\subset U(n)_{\mathrm{reg}}$ into disjoint pieces each inside one $W$ and adding, the identity holds for all Borel $E$, hence for all bounded Borel $f$ in place of
$\mathbf 1_E$. By Step 1 we may replace $T_{\mathrm{reg}}$ by $T$ and $U(n)_{\mathrm{reg}}$ by $U(n)$:
\[\int_{\mathcal O\times T}f\circ\Phi\,\dd|\Phi^*\omega|=n!\int_{U(n)}f\dd\mu.\]

\emph{Step 3: class functions.} If $f$ is a class function, $f(\Phi(gDg^{-1},t))=f(gtg^{-1})=f(t)$. By Lemma~\ref{lem:jacobian} and property (iii),
\[n!\int f\dd\mu=C\int_{\mathcal O}\dd|\sigma|\int_Tf(t)|\Delta(t)|^2\dd\theta.\]
So $\int f\dd\mu=\kappa\int_Tf|\Delta|^2\dd\theta$ with $\kappa>0$ independent of $f$. For $f=1$: expanding $\Delta(t)=\sum_{\pi\in S_n}\operatorname{sgn}(\pi)\prod_kt_k^{\pi(k)-1}$, we get
\[|\Delta(t)|^2=\sum_{\pi,\rho\in S_n}\operatorname{sgn}(\pi)\operatorname{sgn}(\rho)\prod_ke^{i(\pi(k)-\rho(k))\theta_k},\]
and integrating over $[0,2\pi]^n$ kills every term except those with $\pi=\rho$, so $\int_T|\Delta|^2\dd\theta=(2\pi)^nn!$ and $\kappa=\frac1{(2\pi)^nn!}$.
\end{proof}

\begin{think}
Before reading on, look at the factor $|e^{i\theta_j}-e^{i\theta_k}|^2$: it vanishes when two eigenvalues collide, quadratically. Where does the \emph{square} come from in the proof? (Each pair $j<k$ contributed a complex
coordinate $K_{jk}$, that is, two real dimensions.) For real orthogonal or real symmetric matrices, the off-diagonal entries are real, and one expects $|t_j-t_k|$ to the first power. Can you
guess the exponent for quaternionic matrices? (The ``threefold way'' of Dyson, 1962, gives the exponents $1$, $2$ and $4$.) John von Neumann and Eugene Wigner noticed the same count in 1929: real
symmetric matrices with a double eigenvalue form a set of codimension $2$, Hermitian ones of codimension $3$, so a curve of matrices generically never meets them, and eigenvalues ``avoid crossing''.
\end{think}

\begin{remark}
The same proof gives Weyl's formula for every compact connected Lie group $G$ with maximal torus $T$ and Weyl group $W$: for class functions, $\int_Gf\dd\mu=\frac1{|W|}\int_Tf(t)\det\bigl(\mathrm{Ad}(t^{-1})-I\bigr)_{\mathfrak g/\mathfrak t}\dd t$.
For $U(n)$, $|W|=n!$ (Theorem~\ref{thm:torus}) and the determinant is $\prod_{j\neq k}(1-t_k/t_j)=|\Delta(t)|^2$, the $\det A$ of Lemma~\ref{lem:jacobian}. Weyl combined it with the orthogonality of characters to
find the characters of all irreducible representations; Volume~9 tells that story.
\end{remark}

\section{Correlations, and Dyson's formula}

\begin{lemma}[Andreief, 1883]\label{lem:andreief}
Let $\nu$ be a finite measure on a set $X$ and $\phi_1,\dots,\phi_n$, $\psi_1,\dots,\psi_n$ bounded measurable complex functions. Then
\[\frac1{n!}\int_{X^n}\det\bigl(\phi_j(x_k)\bigr)\det\bigl(\psi_j(x_k)\bigr)\dd\nu(x_1)\cdots\dd\nu(x_n)=\det\Bigl(\int_X\phi_j\psi_k\dd\nu\Bigr)_{j,k\le n}.\]
\end{lemma}

\begin{proof}
Expand both determinants and use Fubini: with $M_{ab}=\int\phi_a\psi_b\dd\nu$, the integral is
\[\sum_{\pi,\rho\in S_n}\operatorname{sgn}(\pi)\operatorname{sgn}(\rho)\prod_kM_{\pi(k)\rho(k)}.\] For fixed $\pi$, put $a=\pi(k)$: the product is $\prod_aM_{a,\rho\pi^{-1}(a)}$,
and $\operatorname{sgn}(\pi)\operatorname{sgn}(\rho)=\operatorname{sgn}(\rho\pi^{-1})$; summing over $\rho$ gives $\det M$. The $n!$ choices of $\pi$ give $n!\det M$.
\end{proof}

\begin{theorem}[Heine's identity]\label{thm:heine}
Let $f$ be a bounded Borel function on the unit circle, with Fourier coefficients $\hat f(m)=\int_0^{2\pi}f(e^{i\theta})e^{-im\theta}\frac{\dd\theta}{2\pi}$. If $U$ is Haar distributed on $U(n)$ with
eigenvalues $\lambda_1,\dots,\lambda_n$,
\[\mathbb E\prod_{j=1}^nf(\lambda_j)=\det\bigl(\hat f(j-k)\bigr)_{0\le j,k\le n-1}.\]
\end{theorem}

\begin{proof}
$\prod_jf(\lambda_j)$ is a bounded Borel class function. In Weyl's formula write $|\Delta|^2=\det(e^{ij\theta_k})\overline{\det(e^{ij\theta_k})}$ (rows $j=0,\dots,n-1$) and apply Lemma~\ref{lem:andreief} with
$\nu=\frac{\dd\theta}{2\pi}$, $\phi_j(\theta)=f(e^{i\theta})e^{ij\theta}$ and $\psi_k(\theta)=e^{-ik\theta}$: the result is $\det\bigl(\hat f(k-j)\bigr)$, the determinant of the transpose.
\end{proof}

Matrices with constant diagonals, $(a_{j-k})$, are \emph{Toeplitz matrices}, and Heine's identity turns questions about random unitary matrices into questions about Toeplitz determinants,
a subject with its own long history (Szeg\H o's limit theorem of 1915 and 1952 describes their growth). For instance, the probability that no eigenvalue falls in a given arc is a Toeplitz
determinant (Exercise~\ref{ex:gap}). To see the repulsion of the eigenvalues we need the joint law of a few of them. Let
\[K_n(\theta,\varphi)=\sum_{a=0}^{n-1}e^{ia(\theta-\varphi)},\qquad\text{so}\qquad|K_n(\theta,\varphi)|=\Bigl|\frac{\sin\frac{n(\theta-\varphi)}2}{\sin\frac{\theta-\varphi}2}\Bigr|.\]

\begin{lemma}[Gaudin]\label{lem:gaudin}
Let $\nu$ be a finite measure and $K$ a bounded measurable function with $\int K(x,y)K(y,z)\dd\nu(y)=K(x,z)$ and $\int K(x,x)\dd\nu(x)=c$. Then for $m\ge1$
\[\int\det\bigl(K(x_a,x_b)\bigr)_{a,b\le m}\dd\nu(x_m)=(c-m+1)\det\bigl(K(x_a,x_b)\bigr)_{a,b\le m-1}.\]
\end{lemma}

\begin{proof}
Expand
\[\det\bigl(K(x_a,x_b)\bigr)_{a,b\le m}=\sum_{\pi\in S_m}\operatorname{sgn}(\pi)\prod_aK(x_a,x_{\pi(a)})\]
and integrate in $x_m$. If $\pi(m)=m$, the factor $K(x_m,x_m)$ integrates to $c$, leaving the term of $\pi$
restricted to $\{1,\dots,m-1\}$; these terms give $c\det_{m-1}$. If $\pi(j)=m$ and $\pi(m)=l$ with $j,l<m$, the two factors containing $x_m$ integrate to $K(x_j,x_l)$, leaving the term of the
permutation $\pi'$ of $\{1,\dots,m-1\}$ with $\pi'(j)=l$, which is $\pi$ with $m$ removed from its cycle; removing one element from a cycle changes the sign. Each $\pi'$ arises in this way
from exactly $m-1$ permutations $\pi$ (insert $m$ after any of the $m-1$ elements, in its cycle), so these terms give $-(m-1)\det_{m-1}$.
\end{proof}

\begin{theorem}[correlation functions of random unitary matrices]\label{thm:correlations}
Let $U$ be Haar distributed on $U(n)$, with eigenvalues $e^{i\theta_1},\dots,e^{i\theta_n}$, and let $1\le k\le n$. For every bounded Borel $g$ on $[0,2\pi)^k$,
\[\mathbb E\sum_{j_1,\dots,j_k\text{ distinct}}g(\theta_{j_1},\dots,\theta_{j_k})=\int_{[0,2\pi)^k}g(\varphi)\det\bigl(K_n(\varphi_a,\varphi_b)\bigr)_{a,b\le k}\frac{\dd\varphi_1\cdots\dd\varphi_k}{(2\pi)^k}.\]
\end{theorem}

\begin{proof}
Let $A$ be the matrix $(e^{ij\theta_k})$ with $0\le j<n$ and $1\le k\le n$. Then $(A^*A)_{kl}=\sum_je^{ij(\theta_l-\theta_k)}=K_n(\theta_l,\theta_k)$, so, up to a transpose,
\[|\Delta|^2=|\det A|^2=\det A^*A=\det\bigl(K_n(\theta_a,\theta_b)\bigr)_{a,b\le n}.\] The sum
on the left is a bounded class function of $U$, so Weyl's formula applies. The density $|\Delta|^2$ is symmetric, so each of the $\frac{n!}{(n-k)!}$ terms of the sum contributes the same, and
\[\mathbb E\sum=\frac1{(n-k)!}\int g(\theta_1,\dots,\theta_k)\det\bigl(K_n(\theta_a,\theta_b)\bigr)_{a,b\le n}\frac{\dd\theta}{(2\pi)^n}.\]
The kernel $K_n$ satisfies the hypotheses of Lemma~\ref{lem:gaudin} for $\nu=\frac{\dd\theta}{2\pi}$ with $c=n$, by orthogonality of the $e^{ia\theta}$. Integrating out $\theta_n,\theta_{n-1},\dots,\theta_{k+1}$ produces the factors
$1,2,\dots,n-k$, that is $(n-k)!$.
\end{proof}

For $k=1$ the density is $K_n(\varphi,\varphi)=n$: the eigenvalues are spread evenly on the circle, on average $\frac n{2\pi}$ per unit of angle. For $k=2$ it is $n^2-|K_n(\varphi_1,\varphi_2)|^2$, which vanishes
on the diagonal: two eigenvalues are never close. To see the law at the scale of the mean spacing $\frac{2\pi}n$, measure angles in that unit.

\begin{corollary}[Dyson's pair correlation]\label{cor:dyson}
For $x\in\R$ let $\langle x\rangle\in(-\pi,\pi]$ be the representative of $x$ modulo $2\pi$. For every continuous $h$ with compact support on $\R$,
\[\lim_{n\to\infty}\frac1n\,\mathbb E\sum_{j\neq k}h\Bigl(\frac n{2\pi}\langle\theta_j-\theta_k\rangle\Bigr)=\int_{\R}h(u)\Bigl(1-\Bigl(\frac{\sin\pi u}{\pi u}\Bigr)^2\Bigr)\dd u.\]
For $n$ independent uniform points in place of the eigenvalues, the limit is $\int h(u)\dd u$.
\end{corollary}

\begin{proof}
By Theorem~\ref{thm:correlations} with $k=2$, the expectation is $\int\!\!\int h\bigl(\frac n{2\pi}\langle\varphi_1-\varphi_2\rangle\bigr)\bigl(n^2-|K_n|^2\bigr)\frac{\dd\varphi_1\dd\varphi_2}{(2\pi)^2}$. The integrand depends only on
$x=\langle\varphi_1-\varphi_2\rangle$, and for fixed $\varphi_2$ the map $\varphi_1\mapsto x$ preserves the measure $\frac{\dd\varphi}{2\pi}$ on the circle, so the expectation divided by $n$ is
\[\frac1n\int_{-\pi}^{\pi}h\Bigl(\frac{nx}{2\pi}\Bigr)\Bigl(n^2-\frac{\sin^2\frac{nx}2}{\sin^2\frac x2}\Bigr)\frac{\dd x}{2\pi}=\int_{-n/2}^{n/2}h(u)\Bigl(1-\frac{\sin^2\pi u}{n^2\sin^2\frac{\pi u}n}\Bigr)\dd u,\]
substituting $x=\frac{2\pi u}n$. The integrand lies in $[0,\|h\|_\infty]$ (since $|K_n|\le n$), and $n\sin\frac{\pi u}n\to\pi u$; once $[-\frac n2,\frac n2]$ contains the support of $h$, dominated convergence
gives the limit. For independent uniform points, each of the $n(n-1)$ ordered pairs has $\langle\theta_j-\theta_k\rangle$ uniform on $(-\pi,\pi]$, and the same substitution gives
$\frac{n-1}n\int_{-n/2}^{n/2}h(u)\dd u\to\int h$.
\end{proof}

\begin{figure}[ht]
\centering
\fig{ch08-cue}
\caption{Left: the $24$ eigenvalues of a Haar random matrix of $U(24)$, sampled as in Theorem~\ref{thm:qrhaar}. Middle: $24$ independent uniform points. The eigenvalues are spread more
evenly; the independent points form clumps and leave gaps. Right: the pair correlation $1-(\frac{\sin\pi u}{\pi u})^2$ of Corollary~\ref{cor:dyson}, against the constant $1$ of independent points;
near $u=0$ it behaves like $\frac{\pi^2}3u^2$.}
\label{fig:cue}
\end{figure}

\begin{beauty}
Write the non-trivial zeros of the Riemann zeta function as $\frac12+i\gamma$ (assuming the Riemann hypothesis, so that the $\gamma$ are real), and rescale them by the local mean spacing:
$\tilde\gamma=\frac\gamma{2\pi}\log\frac\gamma{2\pi e}$; by the Riemann--von Mangoldt formula of Volume~5, Chapter~6, the number of zeros up to height $T$ is $\tilde T+O(\log T)$, so the rescaled zeros have, on average, one per unit length. \emph{Montgomery's pair correlation
conjecture} says that, as $T\to\infty$,
\[\frac1{N(T)}\#\bigl\{(\gamma,\gamma'):0<\gamma\neq\gamma'\le T,\ \tilde\gamma-\tilde\gamma'\in[a,b]\bigr\}\longrightarrow\int_a^b\Bigl(1-\Bigl(\frac{\sin\pi u}{\pi u}\Bigr)^2\Bigr)\dd u,\]
with $N(T)$ the number of zeros up to height $T$: exactly the law of Corollary~\ref{cor:dyson}. In 1973, assuming the Riemann hypothesis, Montgomery proved a weighted form of it for the test functions
whose Fourier transform vanishes outside $[-1,1]$; his method is the explicit formula of Volume~5, Chapter~7. In 1987 Andrew Odlyzko computed $10^5$ consecutive zeros near the $10^{12}$-th and found
their pair correlation, and the distribution of their spacings, matching the random matrix curves with startling precision; a numerical calculation repeats a small version of his
experiment. Nobody knows why. The guess of Hilbert and P\'olya, that the $\gamma$ are the eigenvalues of some unknown self-adjoint operator, would explain it, and it is the most beautiful open
door between the subjects of this series. Volume~11 opens it further.
\end{beauty}

\begin{remark}
Dyson introduced the \emph{circular unitary ensemble}, the Haar measure on $U(n)$, in 1962, as a model of a quantum system with no special symmetry, and computed its correlations with the methods that Madan Lal Mehta and Michel Gaudin had developed in 1960 for Gaussian random matrices. Since the Hermitian ensembles need more analysis (Hermite polynomials in
place of $e^{ia\theta}$), the unitary group is the cleanest place to prove it, and the local laws coincide.
\end{remark}

\section{Exercises}

\begin{exercise}[$\star$]\label{ex:axbmodular}
Let $G$ be the group of matrices $h=\begin{pmatrix}a&b\\0&1\end{pmatrix}$ with $a>0$, $b\in\R$ (the maps $x\mapsto ax+b$). (a) Find $\mathfrak g$ and compute $\det\mathrm{Ad}(h^{-1})$. (b) Show that the left
invariant form of Theorem~\ref{thm:haarform} is a multiple of $\frac{\dd a\wedge\dd b}{a^2}$, and find a right invariant form. Compare with Volume~7, Exercise~8.1.
\end{exercise}

\begin{exercise}[$\star\star$]\label{ex:firstcolumn}
Let $U$ be Haar distributed on $U(n)$. (a) Show that the first column of $U$ is uniformly distributed on the unit sphere of $\C^n$, in the sense that it has the law of $z/|z|$ for a standard
complex Gaussian vector $z$. (b) Show that $|U_{11}|^2$ has density $(n-1)(1-x)^{n-2}$ on $[0,1]$. What is the law of $n|U_{11}|^2$ as $n\to\infty$?
\end{exercise}

\begin{exercise}[$\star\star$]\label{ex:su2weyl}
\emph{How far does a random rotation turn?} (a) The eigenvalues of $U\in SU(2)$ are $e^{\pm i\alpha}$ with $\alpha\in[0,\pi]$. Show that for $U$ Haar distributed, $\alpha$ has density $\frac2\pi\sin^2\alpha$. (b) Deduce that
the angle $\phi\in[0,\pi]$ of a Haar random rotation of $SO(3)$ has density $\frac{1-\cos\phi}\pi$, and that its mean is $\frac\pi2+\frac2\pi$, about $126.5^\circ$.
\end{exercise}

\begin{exercise}[$\star\star$]\label{ex:tracemoments}
Let $U$ be Haar distributed on $U(n)$ and $k\ge1$. Show that $\mathbb E\operatorname{tr}U^k=0$ and $\mathbb E|\operatorname{tr}U^k|^2=\min(k,n)$. For $n$ independent uniform points $e^{i\theta_j}$ of the circle,
what is $\mathbb E\bigl|\sum_je^{ik\theta_j}\bigr|^2$? Interpret.
\end{exercise}

\begin{exercise}[$\star\star$]\label{ex:gap}
Let $0<\alpha<\pi$. Show that the probability that a Haar random $U\in U(n)$ has no eigenvalue $e^{i\theta}$ with $|\theta|<\alpha$ is the Toeplitz determinant $\det\bigl(c_{j-k}\bigr)_{0\le j,k<n}$, with
$c_0=1-\frac\alpha\pi$ and $c_m=-\frac{\sin m\alpha}{\pi m}$ for $m\neq0$. Compute it for $n=2$, and check it for $\alpha$ small and for $\alpha\to\pi$.
\end{exercise}

\begin{exercise}[$\star$]\label{ex:repulsion}
For $n=2$, write the joint density of the two eigenangles of a Haar random $U\in U(2)$, and compute the probability that they are at angular distance less than $\varepsilon$. Compare
with two independent uniform points as $\varepsilon\to0$.
\end{exercise}

\section{Hints}

\begin{hint}{ex:axbmodular}
$\mathfrak g$ is the matrices $\begin{pmatrix}x&y\\0&0\end{pmatrix}$. Left invariant $1$-forms are the entries of $h^{-1}\dd h$, right invariant ones those of $\dd h\,h^{-1}$.
\end{hint}

\begin{hint}{ex:firstcolumn}
Use Theorem~\ref{thm:qrhaar}: the first column of $Q$ is $z_1/|z_1|$. For (b), $|z_j|^2$ are independent exponential variables of mean $1$.
\end{hint}

\begin{hint}{ex:su2weyl}
Example~\ref{ex:su2haar} and Exercise~\ref{ex:firstcolumn} with $n=2$: $|a|^2$ is uniform on $[0,1]$, independent of the argument of $a$, so $a$ is uniform in the unit disc. The trace is $2\operatorname{Re}a$.
\end{hint}

\begin{hint}{ex:tracemoments}
$|\operatorname{tr}U^k|^2=n+\sum_{j\neq l}e^{ik(\theta_j-\theta_l)}$; use Theorem~\ref{thm:correlations} with $k=2$ and expand $|K_n|^2$ in exponentials.
\end{hint}

\begin{hint}{ex:gap}
Heine's identity with $f$ the indicator of the complementary arc.
\end{hint}

\begin{hint}{ex:repulsion}
Weyl's formula with $n=2$; $|e^{i\theta_1}-e^{i\theta_2}|^2=2-2\cos(\theta_1-\theta_2)$.
\end{hint}

\section{Solutions}

\begin{solution}{ex:axbmodular}
\idea The modular function measures how conjugation distorts volume in the Lie algebra; for the affine group it stretches the translations.
\writeup (a) The curves $e^{sX}$ in $G$ have $X=\begin{pmatrix}x&y\\0&0\end{pmatrix}$, and every such $X$ works. With $h^{-1}=\begin{pmatrix}1/a&-b/a\\0&1\end{pmatrix}$,
\[h^{-1}Xh=\begin{pmatrix}x&\frac{xb+y}a\\0&0\end{pmatrix},\]
so in the coordinates $(x,y)$, $\mathrm{Ad}(h^{-1})$ has matrix $\begin{pmatrix}1&0\\b/a&1/a\end{pmatrix}$ and determinant $\frac1a$. (b) $h^{-1}\dd h=\begin{pmatrix}\dd a/a&\dd b/a\\0&0\end{pmatrix}$; its entries are the forms $\theta^i$ of the proof
of Theorem~\ref{thm:haarform}, so the left invariant form is a multiple of $\frac{\dd a\wedge\dd b}{a^2}$. By the theorem, $\mu(Eh)=\frac1a\mu(E)$: the group is not unimodular. Similarly $\dd h\,h^{-1}$ has entries
$\frac{\dd a}a$ and $\dd b-\frac ba\dd a$, whose wedge product $\frac{\dd a\wedge\dd b}a$ is right invariant. These are the measures of Volume~7, Exercise~8.1.
\end{solution}

\begin{solution}{ex:firstcolumn}
\idea Gram--Schmidt does nothing to the first column but normalize it, and a Gaussian vector divided by its length is uniform on the sphere.
\writeup (a) By Theorem~\ref{thm:qrhaar} we may take $U=Q$, whose first column is $z_1/|z_1|$, with $z_1$ the first column of $Z$, a standard complex Gaussian vector. (b) $|U_{11}|^2=\frac{E_1}{E_1+S}$ with
$E_j=|z_{j1}|^2$ and $S=E_2+\dots+E_n$. Each $E_j=\frac{X^2+Y^2}2$ is exponential of mean $1$ (its law is $\mathbb P(E_j>s)=e^{-s}$, by polar coordinates in the plane), and they are independent. For $0<x<1$,
\[\mathbb P\Bigl(\frac{E_1}{E_1+S}>x\Bigr)=\mathbb P\Bigl(E_1>\frac{xS}{1-x}\Bigr)=\mathbb E\,e^{-xS/(1-x)}=\Bigl(1+\frac x{1-x}\Bigr)^{-(n-1)}=(1-x)^{n-1},\]
using $\mathbb Ee^{-sE_j}=\frac1{1+s}$. Differentiating, the density is $(n-1)(1-x)^{n-2}$. Then $\mathbb P(n|U_{11}|^2>s)=(1-\frac sn)^{n-1}\to e^{-s}$: an entry of a large random unitary matrix looks like
a complex Gaussian of variance $\frac1n$.
\end{solution}

\begin{solution}{ex:su2weyl}
\idea Under the Haar measure of $SU(2)$, the top-left entry $a$ is a uniform point of the unit disc, and the trace is twice its real part.
\writeup (a) By Example~\ref{ex:su2haar}, $(a,b)=z/|z|$ with $z=(z_1,z_2)$ standard complex Gaussian. By Exercise~\ref{ex:firstcolumn} with $n=2$, $|a|^2$ is uniform on $[0,1]$. The argument of $a$ is that of $z_1$,
uniform on $[0,2\pi)$ and independent of $|z_1|$ and $|z_2|$, hence of $|a|$. So $a=\sqrt V e^{i\psi}$ with $V$, $\psi$ independent and uniform, which is a uniform point of the disc: $\mathbb P(|a|\le r)=r^2$ is the area
fraction. The eigenvalues $e^{\pm i\alpha}$ have sum $2\cos\alpha=\operatorname{tr}U=2\operatorname{Re}a$, and $\operatorname{Re}a$, the abscissa of a uniform point of the disc, has density $\frac2\pi\sqrt{1-x^2}$ on $[-1,1]$.
With $x=\cos\alpha$, $\alpha$ has density $\frac2\pi\sin\alpha\cdot\sin\alpha$ on $[0,\pi]$.
(b) Write $U=e^{\theta X(u)}$ as in Theorem~\ref{thm:spin}; its eigenvalues are $e^{\pm i\theta/2}$, and its image is the rotation by $\theta$ about $u$. So $\theta=2\alpha\in[0,2\pi]$ and the rotation angle is
$\phi=\min(2\alpha,2\pi-2\alpha)$. The image of the Haar measure is the Haar measure of $SO(3)$ (Corollary~\ref{cor:pushhaar}). The two halves $\alpha\le\frac\pi2$ and $\alpha\ge\frac\pi2$ contribute equally (the density
is symmetric about $\frac\pi2$), so $\phi$ has density $2\cdot\frac12\cdot\frac2\pi\sin^2\frac\phi2=\frac{1-\cos\phi}\pi$. Its mean is $\frac1\pi\int_0^\pi\phi(1-\cos\phi)\dd\phi=\frac\pi2-\frac1\pi\bigl[\phi\sin\phi+\cos\phi\bigr]_0^\pi=\frac\pi2+\frac2\pi\approx2.207$, that is
about $126.5^\circ$. Small rotations are rare: the density vanishes like $\frac{\phi^2}{2\pi}$, because the rotations by less than $\phi$ form a ball of radius $\phi$ in the model of Figure~\ref{fig:ballmodel}.
\end{solution}

\begin{solution}{ex:tracemoments}
\idea A Haar random matrix has the low frequencies of its eigenvalues almost frozen: this is the rigidity behind the repulsion.
\writeup By Theorem~\ref{thm:correlations} with $k=1$, $\mathbb E\operatorname{tr}U^k=\int_0^{2\pi}ne^{ik\varphi}\frac{\dd\varphi}{2\pi}=0$. Next, $|\operatorname{tr}U^k|^2=\sum_{j,l}e^{ik(\theta_j-\theta_l)}=n+\sum_{j\neq l}e^{ik(\theta_j-\theta_l)}$, and with $k=2$ in the theorem,
\[\mathbb E|\operatorname{tr}U^k|^2=n+\int\!\!\int e^{ik(\varphi_1-\varphi_2)}\bigl(n^2-|K_n(\varphi_1,\varphi_2)|^2\bigr)\frac{\dd\varphi_1\dd\varphi_2}{(2\pi)^2}.\]
The $n^2$ term integrates to $0$. Since
\[|K_n|^2=\sum_{a,b=0}^{n-1}e^{i(a-b)(\varphi_1-\varphi_2)},\]
the integral of $e^{ik(\varphi_1-\varphi_2)}|K_n|^2$ counts the pairs with $a-b=-k$, of which there are $\max(n-k,0)$. Hence
$\mathbb E|\operatorname{tr}U^k|^2=n-\max(n-k,0)=\min(k,n)$. For independent points, the cross terms have mean $0$ and $\mathbb E|\sum_je^{ik\theta_j}|^2=n$ for every $k\ge1$. The eigenvalues cancel each other much
better at low frequencies: for $k=1$ the fluctuation is $1$ instead of $n$. This is the theorem of Diaconis and Shahshahani (1994) for the second moments.
\end{solution}

\begin{solution}{ex:gap}
\idea ``No eigenvalue in the arc'' is ``$\prod_jf(\lambda_j)=1$'' for an indicator $f$, and Heine's identity turns its probability into a determinant.
\writeup Take $f=\mathbf 1$ on the arc $\alpha\le\theta\le2\pi-\alpha$ and $0$ elsewhere; then $\prod_jf(\lambda_j)$ is the indicator of the event, and Theorem~\ref{thm:heine} gives $\det(\hat f(j-k))$. Here
$\hat f(0)=\frac{2\pi-2\alpha}{2\pi}=1-\frac\alpha\pi$, and for $m\neq0$, $\hat f(m)=\frac1{2\pi}\int_\alpha^{2\pi-\alpha}e^{-im\theta}\dd\theta=-\frac1{2\pi}\int_{-\alpha}^{\alpha}e^{-im\theta}\dd\theta=-\frac{\sin m\alpha}{\pi m}$. For $n=2$ the
probability is $\bigl(1-\frac\alpha\pi\bigr)^2-\frac{\sin^2\alpha}{\pi^2}$. As $\alpha\to\pi$ it tends to $0$, as it should. For $\alpha$ small it is $1-\frac{2\alpha}\pi+\frac{\alpha^2-\sin^2\alpha}{\pi^2}=1-\frac{2\alpha}\pi+O(\alpha^4)$: the
probability of an eigenvalue in the arc is the expected number $2\cdot\frac{2\alpha}{2\pi}$, up to the tiny probability of two eigenvalues there, of order $\alpha^4$; for two independent points that
probability would be $\frac{\alpha^2}{\pi^2}$.
\end{solution}

\begin{solution}{ex:repulsion}
\idea For $n=2$ Weyl's density is $1-\cos(\theta_1-\theta_2)$: zero when the two eigenvalues meet, quadratically.
\writeup By Theorem~\ref{thm:weyl}, the density with respect to $\frac{\dd\theta_1\dd\theta_2}{(2\pi)^2}$ is $\frac12|e^{i\theta_1}-e^{i\theta_2}|^2=1-\cos(\theta_1-\theta_2)$. As in the proof of Corollary~\ref{cor:dyson},
$x=\langle\theta_1-\theta_2\rangle$ has density $\frac{1-\cos x}{2\pi}$ on $(-\pi,\pi]$, so the probability that the eigenvalues are at distance less than $\varepsilon$ is $\frac1{2\pi}\int_{-\varepsilon}^\varepsilon(1-\cos x)\dd x=\frac{\varepsilon-\sin\varepsilon}\pi
\sim\frac{\varepsilon^3}{6\pi}$. For independent uniform points it is $\frac\varepsilon\pi$. At $\varepsilon=0.1$ the ratio is about $600$.
\end{solution}

\section*{Sources and further reading}
\addcontentsline{toc}{section}{Sources and further reading}

The meeting of Montgomery and Dyson: K.~Sabbagh, \emph{Dr.\ Riemann's Zeros} (Atlantic Books, 2002), pp.~134--136, which quotes Montgomery's account; the date, 6 April 1972, is
given by the Institute for Advanced Study for its conference ``50 years of number theory and random matrix theory'' (2022). H.~L.~Montgomery, ``The pair correlation of zeros of the zeta function'',
in \emph{Analytic Number Theory}, Proc.\ Sympos.\ Pure Math.\ 24 (American Mathematical Society, 1973), 181--193. F.~J.~Dyson, ``Statistical theory of the energy levels of complex systems.~I'',
\emph{J.\ Math.\ Phys.} 3 (1962), 140--156, and ``The threefold way'', same volume, 1199--1215; ``Birds and frogs'', \emph{Notices Amer.\ Math.\ Soc.} 56 (2009), 212--223, where Dyson tells the
story from his side. A.~M.~Odlyzko, ``On the distribution of spacings between zeros of the zeta function'', \emph{Math.\ Comp.} 48 (1987), 273--308.

A.~Hurwitz, ``\"Uber die Erzeugung der Invarianten durch Integration'', \emph{Nachr.\ Ges.\ Wiss.\ G\"ottingen, Math.-Phys.\ Kl.} (1897), 71--90; A.~Haar, ``Der Massbegriff in der Theorie der kontinuierlichen
Gruppen'', \emph{Ann.\ of Math.} 34 (1933), 147--169; H.~Weyl, ``Theorie der Darstellung kontinuierlicher halb-einfacher Gruppen durch lineare Transformationen I--III'', \emph{Math.\ Z.} 23 (1925), 271--309, and
24 (1926), 328--376 and 377--395. J.~von Neumann and E.~Wigner, ``\"Uber das Verhalten von Eigenwerten bei adiabatischen Prozessen'', \emph{Phys.\ Z.} 30 (1929), 467--470. C.~Andr\'eief, ``Note sur une relation
entre les int\'egrales d\'efinies des produits des fonctions'', \emph{M\'em.\ Soc.\ Sci.\ Phys.\ Nat.\ Bordeaux} (3) 2 (1883), 1--14.

F.~Mezzadri, ``How to generate random matrices from the classical compact groups'', \emph{Notices Amer.\ Math.\ Soc.} 54 (2007), 592--604; K.~Shoemake, ``Uniform random rotations'', in \emph{Graphics Gems III}
(Academic Press, 1992), 124--132; P.~Diaconis and M.~Shahshahani, ``On the eigenvalues of random matrices'', \emph{J.\ Appl.\ Probab.} 31A (1994), 49--62.

For more: E.~S.~Meckes, \emph{The Random Matrix Theory of the Classical Compact Groups} (Cambridge University Press, 2019), Chapters~1--3, whose proof of Weyl's formula is close to ours;
D.~Bump, \emph{Lie Groups}, 2nd ed. (Springer, 2013), for Weyl's formula on every compact Lie group; M.~L.~Mehta, \emph{Random Matrices}, 3rd ed. (Elsevier, 2004), for Gaudin's lemma; and the survey of
J.~P.~Keating and N.~C.~Snaith, ``Random matrix theory and $\zeta(1/2+it)$'', \emph{Comm.\ Math.\ Phys.} 214 (2000), 57--89.
